\subsection{F01 · Preparación del equipo}\label{fase:F01}
\begin{description}[leftmargin=3.2cm,style=nextline]
\item[1. Herramienta] Windows Terminal + PowerShell 7
\item[2. Terminal] PowerShell (usuario normal)
\item[3. Comandos] \leavevmode
\begin{lstlisting}[style=ps]
winget --version
Set-ExecutionPolicy -Scope CurrentUser RemoteSigned
wsl --status
\end{lstlisting}
\item[4. Explicación] Verifica el gestor de paquetes winget, habilita la ejecución de scripts firmados o locales para el usuario actual y consulta si WSL2 está presente (solo se instala si Docker Desktop lo requiere).
\item[5. Agente] Ninguno (fase manual). Contexto: Equipo con Windows 11 actualizado, 16 GB de RAM recomendados para modelos locales de 3B-8B.
\item[6. Archivos] —
\item[7. Pruebas] winget --version devuelve una versión; Get-ExecutionPolicy -Scope CurrentUser devuelve RemoteSigned
\item[8. Resultado] Terminal lista para instalar herramientas.
\item[9. Errores] winget no se reconoce $\rightarrow$ Instalar o actualizar 'App Installer' desde Microsoft Store.; Los scripts .ps1 se bloquean $\rightarrow$ Repetir Set-ExecutionPolicy en la misma cuenta; desbloquear con Unblock-File.
\item[10. Condición] \textbf{winget operativo y política de ejecución configurada.}
\end{description}

\subsection{F02 · Instalación de herramientas}\label{fase:F02}
\begin{description}[leftmargin=3.2cm,style=nextline]
\item[1. Herramienta] winget, instaladores oficiales de uv y Claude Code, npm
\item[2. Terminal] PowerShell
\item[3. Comandos] \leavevmode
\begin{lstlisting}[style=ps]
.\scripts\instalar_herramientas.ps1
.\scripts\instalar_herramientas.ps1 -Opcionales
.\scripts\verificar_entorno.ps1
\end{lstlisting}
\item[4. Explicación] Instala Git, VS Code, Node.js LTS, Ollama, uv, Claude Code y OpenCode de forma idempotente; con -Opcionales agrega GitHub CLI, MiKTeX y Docker Desktop. El verificador imprime una tabla de versiones.
\item[5. Agente] Ninguno (fase manual). Contexto: Repositorio clonado o carpeta del proyecto descomprimida.
\item[6. Archivos] —
\item[7. Pruebas] verificar\_entorno.ps1 termina con 'Entorno verificado.'
\item[8. Resultado] git, uv, node y ollama disponibles en PATH.
\item[9. Errores] El comando instalado no se reconoce $\rightarrow$ Cerrar y abrir la terminal para recargar PATH.; opencode no se encuentra tras npm install -g $\rightarrow$ Agregar \%APPDATA\%\textbackslash{}npm al PATH del usuario.
\item[10. Condición] \textbf{Todas las herramientas obligatorias en estado OK.}
\end{description}

\subsection{F03 · Inicialización del repositorio}\label{fase:F03}
\begin{description}[leftmargin=3.2cm,style=nextline]
\item[1. Herramienta] Git, GitHub CLI y uv
\item[2. Terminal] PowerShell en la carpeta del proyecto
\item[3. Comandos] \leavevmode
\begin{lstlisting}[style=ps]
uv sync
Copy-Item .env.example .env
gh auth login --hostname github.com --git-protocol https --web
.\scripts\publicar_github.ps1
\end{lstlisting}
\item[4. Explicación] Materializa el entorno virtual desde uv.lock, prepara el archivo de secretos local (excluido por .gitignore) y publica todo el proyecto en https://github.com/JulianReneMunoz/fabrica-nacional-autopartes: el script inicializa Git (rama main), crea el primer commit, crea el repositorio con 'gh repo create' (público por defecto; -Privado para privado) y publica la rama.
\item[5. Agente] Ninguno. Contexto: Cuenta de GitHub JulianReneMunoz autenticada con 'gh auth login'.
\item[6. Archivos] .venv/, .env
\item[7. Pruebas] uv run python -c ''import fabrica\_autopartes; print('ok')''; git status sin cambios pendientes; gh repo view JulianReneMunoz/fabrica-nacional-autopartes --web abre el repositorio; La pestaña Actions muestra el flujo CI en ejecución
\item[8. Resultado] Repositorio versionado y entorno reproducible.
\item[9. Errores] uv sync falla por versión de Python $\rightarrow$ uv python install 3.12 y repetir uv sync.
\item[10. Condición] \textbf{Importación del paquete correcta y primer commit publicado.}
\end{description}

\subsection{F04 · Análisis de requisitos}\label{fase:F04}
\begin{description}[leftmargin=3.2cm,style=nextline]
\item[1. Herramienta] OpenCode (agente plan) o Claude Code (subagente analista-requisitos)
\item[2. Terminal] PowerShell $\rightarrow$ opencode
\item[3. Comandos] \leavevmode
\begin{lstlisting}[style=ps]
opencode
claude
\end{lstlisting}
\item[4. Explicación] Con el agente en modo plan (solo lectura) se extraen requisitos de los PDF del taller y de las clases.
\item[5. Agente] analista-requisitos. Contexto: PDF del Taller 1 y de las clases 1-3 copiados en docs/requisitos/fuentes/ (no versionar si son confidenciales).
\begin{promptbox}Leer AGENTS.md y los PDF de docs/requisitos/fuentes/. Construir docs/requisitos/matriz\_trazabilidad.csv con las columnas id;origen;descripcion;tipo;prioridad;componente;historia;caso\_uso;prueba;estado. Clasificar cada fila como Documentado (citar documento y página), Diseño o Extensión. No inventar requisitos. Registrar en docs/requisitos/analisis\_documentos.md la información ausente, ambigua o contradictoria. Redacción impersonal. No modificar archivos fuera de docs/requisitos/.\end{promptbox}
\item[6. Archivos] docs/requisitos/matriz\_trazabilidad.csv, docs/requisitos/analisis\_documentos.md
\item[7. Pruebas] uv run python scripts/generar\_matriz.py
\item[8. Resultado] Matriz con al menos los 20 requisitos del taller y su origen.
\item[9. Errores] El agente atribuye al taller funciones no pedidas $\rightarrow$ Rechazar el cambio y exigir cita de página.
\item[10. Condición] \textbf{Cada requisito documentado tiene página de origen verificable.}
\end{description}

\subsection{F05 · Configuración de agentes}\label{fase:F05}
\begin{description}[leftmargin=3.2cm,style=nextline]
\item[1. Herramienta] OpenCode y Claude Code
\item[2. Terminal] PowerShell
\item[3. Comandos] \leavevmode
\begin{lstlisting}[style=ps]
opencode auth login
opencode auth list
claude
uv run python scripts/sincronizar_skills.py
\end{lstlisting}
\item[4. Explicación] Autentica proveedores. En OpenCode se usa un proveedor con clave de API o Ollama local; la suscripción Claude Pro/Max solo se usa en Claude Code y claude.ai según los términos vigentes.
\item[5. Agente] Ninguno. Contexto: opencode.json, .claude/settings.json, .claude/agents/ y skills/ del repositorio.
\begin{promptbox}Revisar opencode.json y .claude/settings.json. Confirmar que bash está en 'ask' salvo pruebas y lint, que 'git push' está denegado y que la lectura de .env está bloqueada. Listar los agentes y skills cargados y reportar diferencias con AGENTS.md. No editar archivos.\end{promptbox}
\item[6. Archivos] .opencode/skills/, .claude/skills/
\item[7. Pruebas] En OpenCode: /models muestra el modelo elegido; En Claude Code: /agents lista los 4 subagentes
\item[8. Resultado] Ambas consolas operan con permisos mínimos.
\item[9. Errores] OpenCode no muestra modelos de Anthropic $\rightarrow$ Usar clave de API de la Consola u otro proveedor; no usar credenciales de suscripción.; Modelo local lento o sin herramientas $\rightarrow$ Usar un modelo de código $\geq$14B o un proveedor en la nube para el agente; llama3.2:3b se reserva para la aplicación.
\item[10. Condición] \textbf{Agentes visibles y una tarea de solo lectura ejecutada con éxito.}
\end{description}

\subsection{F06 · Integración MCP}\label{fase:F06}
\begin{description}[leftmargin=3.2cm,style=nextline]
\item[1. Herramienta] Servidor MCP fna-conocimiento, Engram (opcional)
\item[2. Terminal] PowerShell
\item[3. Comandos] \leavevmode
\begin{lstlisting}[style=ps]
uv run fabrica build --fuente semilla
uv run pytest tests/integration/test_mcp.py -q
claude mcp list
engram setup opencode
\end{lstlisting}
\item[4. Explicación] Construye la base para que el servidor tenga datos, prueba el protocolo por stdio y registra servidores.
\item[5. Agente] revisor-calidad. Contexto: .mcp.json y la sección mcp de opencode.json.
\begin{promptbox}Usar la herramienta MCP estadisticas\_base y luego buscar\_fragmentos con la consulta ''sede principal''. Reportar los identificadores devueltos y confirmar que las herramientas son de solo lectura.\end{promptbox}
\item[6. Archivos] —
\item[7. Pruebas] 6 pruebas MCP aprobadas
\item[8. Resultado] El agente consulta la base sin permisos de escritura.
\item[9. Errores] El servidor no inicia $\rightarrow$ Ejecutar 'uv run fabrica-mcp' manualmente y revisar stderr.
\item[10. Condición] \textbf{claude mcp list muestra fna-conocimiento conectado.}
\end{description}

\subsection{F07 · Generación de especificaciones}\label{fase:F07}
\begin{description}[leftmargin=3.2cm,style=nextline]
\item[1. Herramienta] OpenSpec
\item[2. Terminal] PowerShell
\item[3. Comandos] \leavevmode
\begin{lstlisting}[style=ps]
npm install -g @fission-ai/openspec@latest
openspec init --tools claude,opencode
openspec validate --all
\end{lstlisting}
\item[4. Explicación] Inicializa OpenSpec (si no existe) y valida especificaciones y cambios propuestos.
\item[5. Agente] analista-requisitos. Contexto: openspec/config.yaml y matriz de trazabilidad.
\begin{promptbox}/opsx:propose agregar-\textless{}capacidad\textgreater{}. Redactar proposal.md, design.md, tasks.md y el delta de spec con Requirements y Scenarios (WHEN/THEN). Cada tarea termina con un comando de verificación ejecutable.\end{promptbox}
\item[6. Archivos] openspec/changes/\textless{}cambio\textgreater{}/
\item[7. Pruebas] openspec validate --all sin errores
\item[8. Resultado] Especificaciones validadas (evidencia: docs/evidencias/openspec\_validate.txt).
\item[9. Errores] Requirement must have at least one scenario $\rightarrow$ Agregar '\#\#\#\# Scenario:' con viñetas **WHEN**/**THEN**.
\item[10. Condición] \textbf{Validación sin errores.}
\end{description}

\subsection{F08 · Diseño de arquitectura}\label{fase:F08}
\begin{description}[leftmargin=3.2cm,style=nextline]
\item[1. Herramienta] Claude Code (plan) + skill fna-arquitectura
\item[2. Terminal] claude
\item[3. Comandos] \leavevmode
\begin{lstlisting}[style=ps]
uv run python scripts/verificar_capas.py
\end{lstlisting}
\item[4. Explicación] Revisa la dirección de dependencias entre capas y registra decisiones en ADR.
\item[5. Agente] revisor-calidad. Contexto: docs/arquitectura/adr/ y docs/diagramas/.
\begin{promptbox}Aplicar la skill fna-arquitectura sobre src/. Ejecutar scripts/verificar\_capas.py. Si aparece una violación, proponer una inversión de dependencia (protocolo en domain/puertos.py) y el ADR asociado.\end{promptbox}
\item[6. Archivos] docs/arquitectura/adr/ADR-0NN-*.md
\item[7. Pruebas] verificar\_capas.py $\rightarrow$ 0 violaciones
\item[8. Resultado] Arquitectura por capas verificable.
\item[9. Errores] scraping -\textgreater{} pipelines $\rightarrow$ Tipar el rastreador contra el protocolo AlmacenCapturas (hallazgo real corregido).
\item[10. Condición] \textbf{0 violaciones de capas.}
\end{description}

\subsection{F09 · Implementación y ejecución del scraping}\label{fase:F09}
\begin{description}[leftmargin=3.2cm,style=nextline]
\item[1. Herramienta] Claude Code + subagente ingeniero-datos (Sonnet); OpenCode como alternativa
\item[2. Terminal] PowerShell en C:\textbackslash{}fabrica-nacional-autopartes (o la terminal integrada de Claude Code)
\item[3. Comandos] \leavevmode
\begin{lstlisting}[style=ps]
uv run pytest tests/unit/test_scraping.py tests/integration/test_rastreo.py -q
uv run fabrica pipeline --max-paginas 60
uv run fabrica stats
.\scripts\ejecutar_pipeline.ps1 -MaxPaginas 60 -Publicar
\end{lstlisting}
\item[4. Explicación] Implementa y prueba el rastreador con un sitio simulado; luego ejecuta el proceso real desde la consola del equipo con Windows 11: 'uv run fabrica pipeline' rastrea autopartesfanalca.com.co y fanalca.com (Bronze), construye Silver y Gold con '--fuente scraping', valida el corpus y genera las gráficas de cada etapa en docs/graficas. El scraping no se ejecuta dentro de un agente sin aprobación: Claude Code solicita permiso porque el comando accede a la red.
\item[5. Agente] ingeniero-datos. Contexto: Skill fna-scraping-medallion, spec openspec/specs/extraccion-datos.
\begin{promptbox}Implementar en src/fabrica\_autopartes/scraping/ un rastreador BFS que: lea robots.txt antes de cada host, descubra sitemaps, respete Disallow y Crawl-delay, use User-Agent identificable, pausa mínima de 2 s, timeout de 20 s, reintentos con backoff ante 429/5xx, filtre dominios y extensiones no HTML y guarde cada captura en Bronze (HTML + JSONL con url, url\_final, fecha, estado, sha256). Escribir pruebas con tests/fixtures/sitio y la biblioteca responses; ninguna prueba accede a Internet.\end{promptbox}
\item[6. Archivos] src/fabrica\_autopartes/scraping/*.py, src/fabrica\_autopartes/pipelines/bronze.py, tests/integration/test\_rastreo.py
\item[7. Pruebas] Pruebas de scraping en verde; data/bronze/manifiesto.jsonl con capturas reales; data/gold/metricas\_etapas.json con fuente 'scraping'; seis PNG en docs/graficas sin la marca 'semilla'
\item[8. Resultado] Capturas de los sitios oficiales en Bronze, base Gold reconstruida y gráficas por etapa.
\item[9. Errores] HTTP 403 o 429 $\rightarrow$ Aumentar FNA\_PAUSA\_ENTRE\_PETICIONES\_S; nunca rotar User-Agent ni eludir bloqueos.; Texto vacío por contenido JavaScript $\rightarrow$ uv sync --extra dinamico y usar scraping/dinamico.py.
\item[10. Condición] \textbf{Al menos 10 páginas capturadas sin rutas prohibidas.}
\end{description}

\subsection{F10 · Construcción de pipelines}\label{fase:F10}
\begin{description}[leftmargin=3.2cm,style=nextline]
\item[1. Herramienta] Claude Code + subagente ingeniero-datos
\item[2. Terminal] claude
\item[3. Comandos] \leavevmode
\begin{lstlisting}[style=ps]
uv run fabrica build --fuente scraping
uv run fabrica validate
\end{lstlisting}
\item[4. Explicación] Genera Silver (limpieza y deduplicación) y Gold (Markdown, fragmentos y base consolidada).
\item[5. Agente] ingeniero-datos. Contexto: Spec openspec/specs/base-conocimiento.
\begin{promptbox}Revisar data/gold/base\_conocimiento.txt generado desde el scraping. Identificar líneas de menú, avisos legales o textos repetidos que sobrevivieron a la limpieza y ajustar las reglas de pipelines/silver.py (no editar Gold a mano). Agregar una prueba por cada regla nueva.\end{promptbox}
\item[6. Archivos] src/fabrica\_autopartes/pipelines/silver.py, tests/unit/test\_pipelines.py
\item[7. Pruebas] fabrica validate $\rightarrow$ aprobado: true
\item[8. Resultado] Corpus limpio y trazable.
\item[9. Errores] Fragmentos \textgreater{} 600 tokens $\rightarrow$ Reducir FNA\_TAMANO\_FRAGMENTO.
\item[10. Condición] \textbf{Validación aprobada y revisión humana del archivo consolidado.}
\end{description}

\subsection{F11 · Implementación de persistencia}\label{fase:F11}
\begin{description}[leftmargin=3.2cm,style=nextline]
\item[1. Herramienta] Claude Code
\item[2. Terminal] PowerShell
\item[3. Comandos] \leavevmode
\begin{lstlisting}[style=ps]
uv run alembic upgrade head
uv run pytest tests/integration/test_persistencia_y_evaluacion.py -q
\end{lstlisting}
\item[4. Explicación] Aplica las migraciones del esquema de auditoría y verifica el repositorio.
\item[5. Agente] ingeniero-datos. Contexto: ADR-003.
\begin{promptbox}Si se agrega una columna al modelo ORM, generar la migración con 'uv run alembic revision --autogenerate -m ''\textless{}descripcion\textgreater{}''', revisarla manualmente y probar upgrade y downgrade sobre una base temporal.\end{promptbox}
\item[6. Archivos] alembic/versions/*.py
\item[7. Pruebas] test\_migracion\_alembic\_crea\_esquema
\item[8. Resultado] Esquema versionado.
\item[9. Errores] Tabla ya existe $\rightarrow$ Usar una base nueva o 'alembic stamp head' si el esquema se creó con create\_all.
\item[10. Condición] \textbf{Migraciones aplicadas sin errores.}
\end{description}

\subsection{F12 · Construcción de la API}\label{fase:F12}
\begin{description}[leftmargin=3.2cm,style=nextline]
\item[1. Herramienta] Claude Code + skill fna-api
\item[2. Terminal] PowerShell
\item[3. Comandos] \leavevmode
\begin{lstlisting}[style=ps]
.\scripts\ejecutar_api.ps1
uv run fabrica openapi
\end{lstlisting}
\item[4. Explicación] Levanta FastAPI con recarga y exporta el contrato OpenAPI.
\item[5. Agente] ingeniero-llm. Contexto: docs/api/openapi.json, postman/.
\begin{promptbox}Aplicar la skill fna-api para agregar el endpoint solicitado. Contrato con Pydantic, sin lógica de negocio en la ruta, errores como subclases de ErrorFNA, prueba con TestClient y actualización de la colección Postman.\end{promptbox}
\item[6. Archivos] src/fabrica\_autopartes/api/, tests/integration/test\_api.py
\item[7. Pruebas] http://127.0.0.1:8000/docs responde 200; Colección Postman sin fallos
\item[8. Resultado] API documentada y probada.
\item[9. Errores] 503 proveedor\_llm\_no\_disponible $\rightarrow$ Iniciar Ollama (ollama serve) o configurar clave del proveedor.
\item[10. Condición] \textbf{Pruebas de API en verde.}
\end{description}

\subsection{F13 · Desarrollo del frontend}\label{fase:F13}
\begin{description}[leftmargin=3.2cm,style=nextline]
\item[1. Herramienta] Claude Code
\item[2. Terminal] PowerShell
\item[3. Comandos] \leavevmode
\begin{lstlisting}[style=ps]
.\scripts\ejecutar_app.ps1
uv run pytest tests/e2e -q
\end{lstlisting}
\item[4. Explicación] Ejecuta Streamlit y sus pruebas extremo a extremo con AppTest.
\item[5. Agente] ingeniero-llm. Contexto: Spec openspec/specs/interfaz-prueba.
\begin{promptbox}Mejorar frontend/app.py sin romper las pruebas e2e: las pestañas Q\&A, Resumen y FAQ deben quedar visibles; cada respuesta muestra fundamentación, confianza, modo de contexto y fuentes.\end{promptbox}
\item[6. Archivos] frontend/app.py, tests/e2e/test\_streamlit.py
\item[7. Pruebas] 3 pruebas e2e aprobadas
\item[8. Resultado] Interfaz lista para la demostración.
\item[9. Errores] Puerto 8501 ocupado $\rightarrow$ uv run streamlit run frontend/app.py --server.port 8502
\item[10. Condición] \textbf{Demostración manual de las tres funcionalidades.}
\end{description}

\subsection{F14 · Integración de componentes}\label{fase:F14}
\begin{description}[leftmargin=3.2cm,style=nextline]
\item[1. Herramienta] Claude Code + subagente revisor-calidad
\item[2. Terminal] claude
\item[3. Comandos] \leavevmode
\begin{lstlisting}[style=ps]
uv run fabrica ask "¿Dónde queda la sede principal?" --proveedor ollama
\end{lstlisting}
\item[4. Explicación] Prueba el recorrido completo con el LLM real.
\item[5. Agente] revisor-calidad. Contexto: Ollama activo con llama3.2:3b.
\begin{promptbox}Ejecutar las preguntas P01, P05, P20 y P24 de evaluacion/preguntas.yaml por CLI con el proveedor ollama y prompts v1, v2 y v3. Registrar respuestas literales en docs/prompts/bitacora\_experimentos.md.\end{promptbox}
\item[6. Archivos] docs/prompts/bitacora\_experimentos.md
\item[7. Pruebas] P20 devuelve la frase de rechazo con v2 y v3; P24 queda bloqueada
\item[8. Resultado] Integración verificada con LLM real.
\item[9. Errores] JSON inválido en v3 $\rightarrow$ El analizador tolerante extrae el objeto; si persiste, bajar temperatura a 0.
\item[10. Condición] \textbf{Recorrido completo sin errores.}
\end{description}

\subsection{F15 · Pruebas}\label{fase:F15}
\begin{description}[leftmargin=3.2cm,style=nextline]
\item[1. Herramienta] pytest, Ruff, mypy, evaluación
\item[2. Terminal] PowerShell
\item[3. Comandos] \leavevmode
\begin{lstlisting}[style=ps]
.\scripts\verificar_calidad.ps1
uv run fabrica evaluate --proveedor ollama --version-prompt v1
uv run fabrica evaluate --proveedor ollama --version-prompt v3
\end{lstlisting}
\item[4. Explicación] Ejecuta la batería de calidad y la evaluación de 24 preguntas por versión de prompt.
\item[5. Agente] revisor-calidad. Contexto: Skill fna-pruebas.
\begin{promptbox}Aplicar la skill fna-pruebas. Comparar los informes de evaluacion/resultados/ de v1 y v3: exactitud, rechazo correcto y tasa fundamentada. Explicar diferencias con ejemplos de respuestas. No inventar cifras.\end{promptbox}
\item[6. Archivos] docs/evidencias/calidad\_*.txt, evaluacion/resultados/*.md
\item[7. Pruebas] 0 pruebas fallidas; Cobertura $\geq$ 85 \%
\item[8. Resultado] Evidencia cuantitativa para el informe.
\item[9. Errores] Evaluación lenta $\rightarrow$ Usar --modo-contexto seleccion o un modelo más pequeño.
\item[10. Condición] \textbf{Métricas registradas.}
\end{description}

\subsection{F16 · Corrección de defectos}\label{fase:F16}
\begin{description}[leftmargin=3.2cm,style=nextline]
\item[1. Herramienta] Claude Code
\item[2. Terminal] claude
\item[3. Comandos] \leavevmode
\begin{lstlisting}[style=ps]
uv run pytest -q --lf
\end{lstlisting}
\item[4. Explicación] Reejecuta solo las pruebas que fallaron en la última corrida.
\item[5. Agente] ingeniero-llm / ingeniero-datos. Contexto: Salida de la prueba fallida y archivo afectado.
\begin{promptbox}Reproducir el fallo con una prueba mínima, corregir la causa raíz (no la prueba) y ejecutar la batería completa. Documentar el hallazgo en el CHANGELOG del informe.\end{promptbox}
\item[6. Archivos] Según defecto
\item[7. Pruebas] uv run pytest -q
\item[8. Resultado] Defecto corregido con prueba de regresión.
\item[9. Errores] —
\item[10. Condición] \textbf{Batería en verde.}
\end{description}

\subsection{F17 · Despliegue}\label{fase:F17}
\begin{description}[leftmargin=3.2cm,style=nextline]
\item[1. Herramienta] Docker Desktop (opcional) o ejecución local
\item[2. Terminal] PowerShell
\item[3. Comandos] \leavevmode
\begin{lstlisting}[style=ps]
docker compose up --build
docker compose exec ollama ollama pull llama3.2:3b
\end{lstlisting}
\item[4. Explicación] Levanta Ollama, API e interfaz en contenedores.
\item[5. Agente] Ninguno. Contexto: Docker Desktop con backend WSL2.
\item[6. Archivos] —
\item[7. Pruebas] http://localhost:8501 y http://localhost:8000/salud responden
\item[8. Resultado] Servicios en contenedores.
\item[9. Errores] WSL2 no instalado $\rightarrow$ wsl --install en PowerShell como administrador y reiniciar.
\item[10. Condición] \textbf{Servicios saludables.}
\end{description}

\subsection{F18 · v2.0 · Diseño del modelo de componentes y articulaciones}\label{fase:F18}
\begin{description}[leftmargin=3.2cm,style=nextline]
\item[1. Herramienta] Claude Code + subagente arquitecto-mecanico-3d (Opus)
\item[2. Terminal] claude --model opus
\item[3. Comandos] \leavevmode
\begin{lstlisting}[style=ps]
uv run pytest tests/unit/test_simulador.py -q
\end{lstlisting}
\item[4. Explicación] Define piezas, subconjuntos, articulaciones, límites, despiece reversible y orden de desmontaje del freno de disco.
\item[5. Agente] arquitecto-mecanico-3d. Contexto: Skill fna-simulador-3d; docs/prompt\_maestro\_v2.md secciones 30 a 33.
\begin{promptbox}Diseñar en src/fabrica\_autopartes/simulador/ el modelo Pydantic de Pieza, Articulacion (fija, rotacional, prismática, esférica, compuesta), Subconjunto y Ensamblaje, y el catálogo del freno de disco con pinza flotante (cubo, disco, espárragos, soporte, cuerpo de pinza, pistón, pastillas, pasadores y tornillos). Toda dimensión se marca 'demostrativo'. Definir límites, retirar\_antes sin ciclos y direcciones de despiece. No redactar procedimientos de reparación ni de seguridad.\end{promptbox}
\item[6. Archivos] src/fabrica\_autopartes/simulador/*.py, tests/unit/test\_simulador.py
\item[7. Pruebas] Jerarquía, identificadores, transformaciones, límites y reversibilidad en verde
\item[8. Resultado] Modelo mecánico validado y documentado.
\item[9. Errores] Ciclo en retirar\_antes $\rightarrow$ Revisar dependencias; la secuencia usa orden topológico.
\item[10. Condición] \textbf{Revisión del revisor-calidad (Opus) sin hallazgos bloqueantes.}
\end{description}

\subsection{F19 · v2.0 · Generación procedural del modelo 3D}\label{fase:F19}
\begin{description}[leftmargin=3.2cm,style=nextline]
\item[1. Herramienta] Claude Code + subagente ingeniero-3d-frontend (Sonnet)
\item[2. Terminal] PowerShell
\item[3. Comandos] \leavevmode
\begin{lstlisting}[style=ps]
uv run fabrica simulador generar
uv run fabrica simulador validar
uv run pytest tests/integration/test_generacion_3d.py -q
\end{lstlisting}
\item[4. Explicación] Genera web/public/modelos/freno\_disco.glb y su manifiesto con trimesh y manifold3d desde la CLI.
\item[5. Agente] ingeniero-3d-frontend. Contexto: ADR del motor geométrico (trimesh frente a Blender CLI y CadQuery).
\begin{promptbox}Implementar simulador/generacion.py: una malla cerrada por pieza según su primitiva, material PBR, nodo GLB con el id de la pieza y manifiesto con sha256, volumen y límites. La salida debe ser reproducible (mismo hash con los mismos parámetros).\end{promptbox}
\item[6. Archivos] src/fabrica\_autopartes/simulador/generacion.py, web/public/modelos/freno\_disco.glb
\item[7. Pruebas] Mallas cerradas; GLB con 15 nodos; hash reproducible
\item[8. Resultado] Modelo GLB demostrativo y trazable.
\item[9. Errores] Malla no cerrada $\rightarrow$ Revisar operaciones booleanas de manifold3d y parámetros positivos.
\item[10. Condición] \textbf{simulador validar $\rightarrow$ aprobado.}
\end{description}

\subsection{F20 · v2.0 · API del simulador, RAG y persistencia en PostgreSQL}\label{fase:F20}
\begin{description}[leftmargin=3.2cm,style=nextline]
\item[1. Herramienta] Claude Code + subagentes ingeniero-3d-frontend e ingeniero-llm (Sonnet)
\item[2. Terminal] PowerShell
\item[3. Comandos] \leavevmode
\begin{lstlisting}[style=ps]
docker compose up -d postgres
uv run alembic upgrade head
uv run pytest tests/integration/test_api_simulador.py tests/integration/test_postgresql.py -q
\end{lstlisting}
\item[4. Explicación] Expone /api/v1/simulador (cinemática, despiece, cálculo, preguntas, estados y eventos) y /api/v1/maquinas.
\item[5. Agente] ingeniero-3d-frontend. Contexto: docs/prompt\_maestro\_v2.md secciones 36 a 38.
\begin{promptbox}Implementar las rutas del simulador con validación Pydantic; rechazar con 422 identificadores inexistentes, valores no finitos o fuera de límites. En /preguntar separar información documentada (fragmentos con URL) de datos del modelo y rechazar preguntas de reparación o seguridad.\end{promptbox}
\item[6. Archivos] src/fabrica\_autopartes/api/routes/simulador.py, src/fabrica\_autopartes/services/simulador.py, alembic/versions/0002\_simulador.py
\item[7. Pruebas] Pruebas de API y PostgreSQL en verde
\item[8. Resultado] API documentada en /docs.
\item[9. Errores] connection refused en PostgreSQL $\rightarrow$ docker compose ps; verificar FNA\_URL\_BASE\_DATOS.
\item[10. Condición] \textbf{Pruebas de API en verde.}
\end{description}

\subsection{F21 · v2.0 · Visor 3D y biblioteca de máquinas (Anatomy)}\label{fase:F21}
\begin{description}[leftmargin=3.2cm,style=nextline]
\item[1. Herramienta] Claude Code + subagente ingeniero-3d-frontend (Sonnet) + skill Anatomy
\item[2. Terminal] PowerShell en web/
\item[3. Comandos] \leavevmode
\begin{lstlisting}[style=ps]
npx skills add wheresryan22/anatomy -g -a claude-code
npm ci
npm run typecheck
npm test
npm run build
\end{lstlisting}
\item[4. Explicación] Construye el visor React Three Fiber y la biblioteca de máquinas con datos del scraping.
\item[5. Agente] ingeniero-3d-frontend. Contexto: Skills fna-maquinas-anatomia y /anatomy; secciones 35 y 39 del prompt maestro v2.
\begin{promptbox}Implementar en web/ la biblioteca (tarjetas con pregunta, una línea, relación, datos de empresa con URL y botón Abrir máquina) y las máquinas freno 3D, prensa y línea. Cada máquina: 2–4 controles, correr/pausar/reiniciar, cámara lenta, cifras junto a las piezas, idea principal, revisión crítica, recorrido ?tour=1 menor a 30 s y modo experimento. El freno incluye órbita, zoom, selección, árbol, ficha, despiece/reensamblaje, visibilidad, transparencia, sección, reinicio de cámara, PNG y lecturas.\end{promptbox}
\item[6. Archivos] web/src/**/*.tsx, web/src/modelos/*.ts, web/src/\_\_tests\_\_/*.ts*
\item[7. Pruebas] Vitest en verde (réplica TS comparada con Python); npm run build sin errores
\item[8. Resultado] Visor servido por FastAPI en http://127.0.0.1:8000/app/.
\item[9. Errores] Pantalla en blanco en /app $\rightarrow$ Ejecutar npm run build; la API monta web/dist solo si existe.
\item[10. Condición] \textbf{Pruebas del visor en verde.}
\end{description}

\subsection{F22 · v2.0 · Pruebas de extremo a extremo y videos del recorrido}\label{fase:F22}
\begin{description}[leftmargin=3.2cm,style=nextline]
\item[1. Herramienta] Playwright con Chrome del sistema + FFmpeg
\item[2. Terminal] PowerShell en web/
\item[3. Comandos] \leavevmode
\begin{lstlisting}[style=ps]
npx playwright test --project=funcional
npx playwright test --project=recorrido
\end{lstlisting}
\item[4. Explicación] Verifica biblioteca, freno 3D, prensa y línea; mide FPS; graba los recorridos en docs/videos/*.mp4.
\item[5. Agente] revisor-calidad. Contexto: web/e2e/.
\begin{promptbox}Ejecutar las pruebas de extremo a extremo y revisar las capturas de docs/capturas/v2 y los videos de docs/videos. Reportar errores de consola, textos superpuestos y cualquier cifra sin procedencia.\end{promptbox}
\item[6. Archivos] docs/capturas/v2/*.png, docs/videos/*.mp4
\item[7. Pruebas] 7 pruebas en verde; videos de menos de 30 s
\item[8. Resultado] Evidencias visuales para la presentación.
\item[9. Errores] ffmpeg no se reconoce $\rightarrow$ winget install --id Gyan.FFmpeg -e y reabrir la terminal.
\item[10. Condición] \textbf{Sin hallazgos bloqueantes.}
\end{description}

\subsection{F23 · v2.0 · Integración en Streamlit}\label{fase:F23}
\begin{description}[leftmargin=3.2cm,style=nextline]
\item[1. Herramienta] Claude Code (Sonnet)
\item[2. Terminal] PowerShell
\item[3. Comandos] \leavevmode
\begin{lstlisting}[style=ps]
.\scripts\ejecutar_api.ps1
uv run streamlit run frontend/app.py
uv run pytest tests/e2e/test_streamlit.py -q
\end{lstlisting}
\item[4. Explicación] Streamlit sigue como interfaz principal: pestañas Proceso de datos (gráficas por etapa) y Simulador 3D (enlace e iframe).
\item[5. Agente] ingeniero-3d-frontend. Contexto: frontend/app.py.
\begin{promptbox}Agregar las pestañas Proceso de datos y Simulador 3D. La primera muestra el flujo, las métricas y las seis gráficas con su explicación; la segunda enlaza e incrusta el visor y lista las máquinas con sus datos y URL.\end{promptbox}
\item[6. Archivos] frontend/app.py, tests/e2e/test\_streamlit.py
\item[7. Pruebas] AppTest sin excepciones
\item[8. Resultado] Interfaz empresarial integrada.
\item[9. Errores] El iframe no carga $\rightarrow$ Iniciar la API y compilar web/; revisar FNA\_URL\_VISOR.
\item[10. Condición] \textbf{Pestañas operativas.}
\end{description}

\subsection{F24 · Documentación final}\label{fase:F24}
\begin{description}[leftmargin=3.2cm,style=nextline]
\item[1. Herramienta] MiKTeX/TeX Live + subagente documentador (Sonnet) + skill fna-documentacion
\item[2. Terminal] PowerShell en docs/latex
\item[3. Comandos] \leavevmode
\begin{lstlisting}[style=ps]
uv run python scripts/generar_matriz.py
uv run python scripts/generar_guia.py
latexmk -pdf -interaction=nonstopmode main.tex
\end{lstlisting}
\item[4. Explicación] Regenera tablas desde las fuentes y compila el manual.
\item[5. Agente] documentador. Contexto: docs/latex/.
\begin{promptbox}Aplicar la skill fna-documentacion. Actualizar el capítulo de resultados únicamente con cifras de evaluacion/resultados/. Mantener redacción impersonal y verificar que no queden referencias sin resolver.\end{promptbox}
\item[6. Archivos] docs/latex/main.pdf, README.md
\item[7. Pruebas] El PDF compila sin errores
\item[8. Resultado] Manual y README listos para la entrega.
\item[9. Errores] Paquete LaTeX faltante $\rightarrow$ MiKTeX instala paquetes bajo demanda; habilitar la instalación automática.
\item[10. Condición] \textbf{Entrega completa.}
\end{description}
